% ============================================================
% CHAPTER 6: PROJECT GOALS
% MCDA 5585 / 5586 Major Project Report
% Bhavik Kantilal Bhagat | A00494758 | Saint Mary's University
% ============================================================

\chapter{Project Goals}
\label{chap:project_goals}

I set the following engineering objectives at the outset of the co-op placement,
guided by the two motivating incidents and the requirements sessions with
Mr.\ Deotale and Ms.\ Labos. Each objective maps directly to one of the six work streams
described in Chapter~\ref{chap:project}.


\FloatBarrier

\section{Project Goals}
\label{sec:project_goals}


The primary goals of the co-op placement project were as follows:

\begin{itemize}

    \item \textbf{Eliminate reactive-only incident detection.} Replace the
          dependency on business-team escalation as the primary detection
          mechanism with proactive CloudWatch alarms that fire within seconds
          of a threshold being crossed, directly addressing the failure modes
          exposed by the CAIS-Pricing Silent Outage and the Meridian Incident.

    \item \textbf{Build a centralized observability platform for the IA-IT-SRE
          Infrastructure.} Provide a single-pane-of-glass view across all six
          platforms and all AWS service categories used by the IA infrastructure, enabling the SRE team
          to assess platform health, drill down to individual AWS services, and
          respond to incidents without navigating six separate CloudWatch consoles.

    \item \textbf{Automate AWS service discovery and inventory.} Replace manual
          Confluence-maintained Lambda lists with a programmatic, always-current
          inventory API that enumerates all deployed AWS resources per platform
          and exposes them via REST endpoints for use by the dashboard and other
          downstream tooling.

    \item \textbf{Evaluate AWS-native observability as a Dynatrace alternative.}
          Conduct a structured feature-parity analysis comparing the AWS-native
          stack (CloudWatch, Grafana AMG, X-Ray, Application Signals) against
          Dynatrace capabilities, quantifying the achievable parity percentage
          and the incremental monthly cost at each tier.

    \item \textbf{Deliver Infrastructure-as-Code for all monitoring components.}
          Ensure all alarms, dashboards, IAM roles, and supporting infrastructure
          are provisioned via parameterized CloudFormation templates and AWS SAM,
          making the platform reproducible, version-controlled, and
          environment-portable.

    \item \textbf{Establish engineering quality standards.} Produce a comprehensive test suite
          for the inventory API, document all
          architectural decisions and limitations, and maintain a Confluence
          knowledge base for the Lambda inventory, runbooks, and incident
          post-mortems.

\end{itemize}


\FloatBarrier

\section{Deliverables and Technical Scope}
\label{sec:deliverables_scope}


This section documents the technical scope and stack for each primary
deliverable. Detailed implementation decisions are covered in
Chapter~\ref{chap:implementation}.


\subsection{IA Infrastructure Monitoring Platform}

\textbf{Current State:} 
No CloudWatch dashboards existed for any of the six
IA platforms. Lambda function health was monitored only through the AWS
Console on an ad-hoc basis. No tag-based resource discovery was in place
and the Lambda inventory was maintained manually on Confluence.

\textbf{Scope:} 
Build unified CloudWatch dashboards covering all six
platforms --- \texttt{automationhub}, \texttt{bpm}, \texttt{citcoworks},
\texttt{ia-bpo}, \texttt{meridian}, and \texttt{rpacfs1} --- across Lambda,
EC2, ECS, EKS, SQS, and supporting AWS services. Implement a tag-based
resource discovery pipeline using budget\_code as the primary filter and
AppName as the mandatory fallback discriminator for the
\texttt{meridian}/\texttt{citcoworks} shared budget\_code ambiguity.

\textbf{Deliverables:} 
CloudWatch dashboards for all six platforms; a
resilient three-tier SQS registry pipeline (live API → S3 cache → bundled
inventory); a Python-based dashboard builder framework deployed via
CloudFormation.


\subsection{Centralized Grafana SRE Dashboard}

\textbf{Current State:} 
No centralized observability workspace existed.
Platform health required navigating six separate CloudWatch consoles.
No Kafka consumer lag alarm existed for the Meridian platform, which was
the direct cause of the April~2026 incident going undetected for over
three hours.

\textbf{Scope:} 
Build a centralized Amazon Managed Grafana (AMG) workspace
providing a unified cross-platform view across all six platforms and all
monitored AWS service categories. Implement proactive CloudWatch alarms
across all services, an SNS notification pipeline to the SRE team, and
Grafana template variables backed by the Inventory API for dynamic
resource filtering.

\textbf{Deliverables:} 
Grafana AMG workspace with 13+ dashboard rows
covering all service categories; CloudWatch alarms deployed across all
six platforms; SNS alarm notification pipeline; Infrastructure-as-Code
(AWS SAM / CloudFormation) for reproducible deployment.


\subsection{IA Infrastructure Inventory API}

\textbf{Current State:} 
No programmatic AWS resource inventory existed.
Lambda function lists were manually maintained on Confluence and became
stale quickly. CloudWatch Metric Insights SQL lacked support for
pattern-based dimension filtering, making dynamic Grafana template variable
population impractical.

\textbf{Scope:} 
Build a Flask REST API that dynamically discovers all
deployed AWS resources per platform across all six IA platforms. Implement
a discoverer registry pattern for extensibility, a two-stage tag filter
to resolve the \texttt{meridian}/\texttt{citcoworks} budget\_code ambiguity,
and an S3-backed cache layer to keep enumeration latency within the
API Gateway 30-second timeout.

\textbf{Deliverables:} 
REST endpoints for Lambda, ECS, EC2, SQS, S3, and
health resources; per-platform filtering support; S3 cache with
CloudWatch metrics publishing; Lambda deployment via API Gateway.


\subsection{RPA Production Support}

\textbf{Current State:} 
The IA-SRE team operates a Follow the Sun support
model with shifts across Manila (APAC), Dublin (EMEA), and Halifax (NA).
The NA slot (11~AM--7~PM~EST) rotated weekly among the four NA-based
SRE team members.

\textbf{Scope:} 
Cover the North America RPA production support rotation,
handling incidents escalated through the SD Portal, IA\_INCIDENTS
distribution list, and Teams channels. Triage, investigate, and resolve
or escalate RPA process failures within the defined SLA, and maintain
daily Confluence shift handover logs.

\textbf{Deliverables:} 
Five completed support rotations (W8, W12, W16,
W20, W24); incident triage and resolution for all escalated tickets;
Confluence handover logs maintained for each shift day.


\subsection{RPA LLM Agent --- Ask Citco RPA}

\textbf{Current State:} 
The \textbf{Ask Citco RPA} chatbot was a
production RAG pipeline using Amazon Bedrock (Claude), OpenSearch
Serverless (1024-dim FAISS/HNSW vectors), and 11 Lambda functions
ingesting from SharePoint, Databricks, and the BPM API. The system
had multiple known reliability issues including a recurring KB sync
timeout and several production bugs identified via CloudWatch log analysis.

\textbf{Scope:} 
Resolve production reliability issues, enrich the
Knowledge Base, and improve code quality across all Lambda functions
per the IA Best Practice checklist. Evaluate and propose a Hybrid
Graph-RAG architecture extension using Amazon Neptune.

\textbf{Deliverables:} 
Production bug fixes deployed; NavChecklist PDFs
added to Knowledge Base; modular code structure with structured logging,
Pydantic validation, tenacity retry, kill switch, and X-Ray tracing;
comprehensive unit test suite; architecture proposal for Graph-RAG extension.


\subsection{Citco Enterprise Custom UI Dashboards}

\textbf{Current State:} 
The Grafana AMG dashboard provided approximately
60\% Dynatrace feature parity but was constrained by fixed panel types,
no drag-and-drop layout, and limited custom visualization options. A
Citco-owned enterprise UI was identified as the next step to close
the remaining observability gap.

\textbf{Scope:} 
Extend the existing \texttt{ia-config-portal}
(React~19/TypeScript, ECS Fargate, built by Ms.\ Venkataraman and Mr.\ Vadla)
with two new dashboard pages --- an Infrastructure Inventory dashboard
and an Infrastructure Monitoring dashboard ---
using a config-driven panel schema and CloudWatch data integration.

\textbf{Deliverables:} 
Architectural implementation plan covering gap
analysis, technology stack decisions, API design, and task breakdown;
initial frontend scaffolding with AG~Grid resource tables and
config-driven panel components; continuing in Sep--Dec~2026 extension.


\FloatBarrier

\section{Acceptance Criteria}
\label{sec:acceptance_criteria}


The following criteria defined whether each deliverable met its stated
objectives.

\subsection{IA Infrastructure Monitoring Platform}

\begin{itemize}
    \item \textbf{Platform coverage}: CloudWatch dashboards exist for
          all six platforms covering the key AWS services (Lambda, ECS,
          EC2, SQS, EKS) with tag-based resource filtering.

    \item \textbf{Resilient SQS monitoring}: The SQS registry pipeline
          remains functional during AWS service degradation via the
          three-tier fallback model.
\end{itemize}


\subsection{SRE Observability Platform}

\begin{itemize}
    \item \textbf{Proactive incident detection}: The team can be alerted
          to Kafka consumer lag, Lambda error rate spikes, ECS CPU
          saturation, and SQS queue backlogs before business teams
          escalate --- directly addressing the failure modes exposed by
          the CAIS-Pricing Silent Outage and the Meridian Incident.

    \item \textbf{Cross-platform visibility}: The SRE team can assess
          the health of all six platforms and all monitored AWS services
          from a single Grafana workspace without navigating individual
          CloudWatch consoles.

    \item \textbf{Alarm coverage}: CloudWatch alarms are deployed and
          active across all monitored AWS services on all six platforms,
          with SNS notifications routed to the SRE team.

    \item \textbf{Container-level telemetry}: Task-level CPU and memory
          metrics are available for the Meridian ECS clusters, enabling
          resource saturation to be detected before service degradation
          occurs.

    \item \textbf{Reproducible infrastructure}: All monitoring
          infrastructure --- Grafana workspace, alarms, IAM roles, S3
          storage --- is provisioned as Infrastructure-as-Code and
          can be redeployed consistently across environments.
\end{itemize}


\subsection{IA Infrastructure Inventory API}

\begin{itemize}
    \item \textbf{Dynamic discovery}: The API returns a current,
          accurate list of deployed AWS resources per platform at query
          time, enabling Grafana template variables to reflect the live
          fleet without manual updates.

    \item \textbf{Tag disambiguation}: Resources from the
          \texttt{meridian} and \texttt{citcoworks} platforms are
          correctly separated despite sharing the same budget\_code
          tag value.

    \item \textbf{Extensibility}: New platforms or AWS service types
          can be onboarded by implementing a single discoverer class
          with no changes to the API routing or test infrastructure.
\end{itemize}

\subsection{RPA Production Support}

\begin{itemize}
    \item \textbf{Incident response}: Production incidents are triaged,
          root-caused, and either resolved within the support shift or
          escalated with a documented RCA and proposed fix.

    \item \textbf{SLA compliance}: Tickets are acknowledged within the
          L2 (Innovation IT) response SLA and carried over with complete
          handover notes per the IA Support Model. L3 (SRE) is escalated
          for incidents requiring infrastructure access or code-level fixes.
\end{itemize}


\subsection{RPA LLM Agent}

\begin{itemize}
    \item \textbf{KB sync reliability}: The \texttt{IRCOEKBExtractLambda}
          completes within the 300-second timeout on every daily run,
          with steady-state duration under 60 seconds for unchanged
          content.

    \item \textbf{Production stability}: All identified bugs (OpenSearch
          silent delete, fire-and-forget pattern, Athena race condition,
          SSL verification, encryption inconsistency) are resolved and
          deployed.

    \item \textbf{KB content}: NavChecklist documents are indexed and
          retrievable via the chatbot.
\end{itemize}


\subsection{Citco Enterprise Custom UI Dashboards}

\begin{itemize}
    \item \textbf{Architecture plan}: A comprehensive implementation plan
          covering gap analysis, technology decisions, API design, and
          task breakdown is produced before implementation begins.

    \item \textbf{Backend API design}: New Lambda handlers designed
          with request/response schemas and IAM permissions documented.

    \item \textbf{Component specifications}: SRE-specific UI components
          (honeycomb, service map, trace waterfall) specified with library
          selections and integration patterns.
\end{itemize}



% NOTE: Team composition and mandate are documented in Chapter 1
% (About Company & Organization, Section 2.2 Team Structure).


% Keep all floats within this chapter
\FloatBarrier
